\bookchapter{Chapitre 69 : Mémo et ressources}

\chapterintro{Ce chapitre est un mémo qui se suffit à lui-même. Quand tu ne te rappelles plus d'une écriture ou qu'il te faut décider vite, ouvre-le directement. Traite-le comme des « cartes de poche » : mémo du code courant, questions fréquentes, ressources d'apprentissage — sans en sauter la moindre page.}

\booksection{La fiche mémo du code courant}

\begin{center}
\begin{tabularx}{\linewidth}{|X|X|}
\hline
\textbf{Ce que tu veux faire} & \textbf{Code clé} \tabularnewline
\hline
État privé de la page & @State private var x = valeurInitiale \tabularnewline
Logique entre les pages & \texttt{@Observable class VM \{ var x = 0 \}} \tabularnewline
Disposition verticale / horizontale & \texttt{VStack \{ \}}/\texttt{HStack \{ \}} \tabularnewline
Espace élastique & \texttt{Spacer()} \tabularnewline
Liste & List \{ ForEach(tableau) \{ ... \} \} \tabularnewline
Saut d'une page à l'autre & \texttt{NavigationStack + NavigationLink} \tabularnewline
Fenêtre en bas d'écran & .sheet(isPresented: \$show) \{ contenu \} \tabularnewline
Fermer la page courante & \texttt{@Environment(\textbackslash{}.\allowbreak dismiss) var dismiss} \tabularnewline
Récupérer des données en ligne & \texttt{try await URLSession.\allowbreak shared.\allowbreak data(from: url)} \tabularnewline
Analyse JSON & \texttt{struct X: Codable}+\texttt{JSONDecoder().\allowbreak decode} \tabularnewline
Stocker une petite configuration & UserDefaults.standard.set(valeur, forKey: "cle") \tabularnewline
Stocker des données structurées & \texttt{@Model}+\texttt{modelContainer}+\texttt{@Query} \tabularnewline
Adaptation au mode sombre & Utiliser les couleurs système (.primary / Color(.systemGray6)), ne rien coder en dur \tabularnewline
Lancer les tests & CmdU \tabularnewline
\hline
\end{tabularx}
\end{center}

\begin{center}\small\color{gray}Aide-mémoire du développement iOS (édition d'août 2026)\end{center}

\booksection{II. Questions fréquentes (FAQ)}

Q1 : Je n'ai pas de Mac, puis-je apprendre le développement iOS ?

Tu peux apprendre le langage, mais pas faire un développement complet. Swift Playgrounds sur iPad/iPhone permet de couvrir tout le contenu du chapitre 61 ; mais Xcode est exclusif à macOS et il te faudra au final un Mac. Avec un budget limité, tu peux d'abord louer un Mac dans le cloud ou acheter un Mac mini d'occasion pour dépanner, puis décider de ton investissement une fois le langage maîtrisé.

Q2 : Combien coûte l'apprentissage du développement iOS ?

Phase d'apprentissage : 0 yuan (Xcode gratuit + tests sur appareil réel avec un identifiant Apple gratuit). Phase de publication : compte développeur à 99 \$ par an (environ 700 yuans). Ajoute l'investissement matériel d'un Mac, et c'est tout le coût. Au regard du retour — même un petit utilitaire sur l'App Store trouve des utilisateurs payants — le calcul est rentable.

Q3 : SwiftUI ou UIKit, lequel apprendre en premier ?

SwiftUI. En 2026, les nouveaux projets, tutoriels et postes sont dominés par SwiftUI, et ce livre l'utilise du début à la fin. UIKit est le framework de la génération précédente, encore très utilisé dans d'anciens projets — une fois SwiftUI bien maîtrisé, complète avec UIKit selon tes besoins (Apple propose des tutoriels d'interopérabilité), tu t'y mettras très vite. N'inverse pas l'ordre.

Q4 : Est-ce facile de trouver du travail en développement iOS ?

L'état du marché en 2026 : les postes juniors qui se contentent de « dessiner des interfaces » sont très disputés, mais les développeurs « à l'aise en Swift + SwiftUI + avec une app publiée » restent très recherchés. Le salaire médian des postes iOS figure depuis toujours dans le haut du classement des différentes technos. Pendant ton apprentissage, ton meilleur CV, c'est l'app que tu as faite toi-même et disponible sur la boutique.

Q5 : Que faire si l'app est refusée lors de la validation ?

Pas de panique, c'est tout à fait normal — une première soumission a de bonnes chances d'être refusée. Apple motive ses refus de façon très précise (quelle page, quelle règle) ; corrige point par point, réponds avec tes explications dans App Store Connect, puis renvoie. Les trois grandes causes de refus : l'usage des permissions mal expliqué, des bugs d'interface, des métadonnées (captures/description) qui ne correspondent pas au produit. Vérifie ces trois points à l'avance et tu économiseras un tour d'attente.

\booksection{III. Aide-mémoire des raccourcis Xcode les plus utilisés}

\begin{center}
\begin{tabularx}{\linewidth}{|X|X|}
\hline
\textbf{Action} & \textbf{Raccourci} \tabularnewline
\hline
Exécuter / Tester & CmdR / CmdU \tabularnewline
Ouvrir l'aperçu du canvas & OptionCmd \tabularnewline
Ouvrir rapidement un fichier & CmdO \tabularnewline
Recherche globale & CmdF \tabularnewline
Réindenter le code & CtrlI \tabularnewline
Commenter / Décommenter & Cmd/ \tabularnewline
Nettoyer le cache de build & CmdK \tabularnewline
Aller à la définition & Cmd + clic \tabularnewline
Renommer (renommage global sécurisé) & Clic droit Refactor Rename \tabularnewline
\hline
\end{tabularx}
\end{center}

\begin{center}\small\color{gray}Aide-mémoire des raccourcis Xcode les plus utilisés\end{center}

\begin{tcbinfo}{Petit conseil : en cas de compilation mystérieuse, vide d'abord le cache}
Les événements paranormaux du type « hier ça tournait, aujourd'hui plus rien » viennent dans six à sept cas sur dix d'un cache de build corrompu. Un CmdK pour nettoyer puis recompiler guérit la grande majorité de ces bizarreries. Si ça ne suffit pas, redémarre Xcode (voire le Mac entier), et en dernier recours « supprime le dossier DerivedData » — c'est le rituel en trois temps que se transmettent les développeurs iOS.

\end{tcbinfo}

\booksection{IV. Ressources d'apprentissage recommandées}

\booksection{VI. Ressources d'apprentissage recommandées}

\begin{center}
\begin{tabularx}{\linewidth}{|X|X|X|}
\hline
\textbf{Type} & \textbf{Ressource} & \textbf{Pourquoi c'est recommandé} \tabularnewline
\hline
Officiel & developer.apple.com & Documentation officielle d'Apple, où se trouvent les tutoriels SwiftUI/Swift : la référence numéro un \tabularnewline
Officiel & Les SwiftUI Tutorials d'Apple & Cours officiel guidé pas à pas, convivialité maximale pour les débutants \tabularnewline
Officiel & Site officiel de Swift (swift.org) & Documentation officielle du langage et bac à sable en ligne \tabularnewline
Communauté & Hacking with Swift & La bibliothèque de tutoriels gratuits de Paul Hudson, dont la qualité est unanimement reconnue parmi les toutes meilleures \tabularnewline
Communauté & Stack Overflow (tag swiftui) & Le premier réflexe pour demander de l'aide face à une erreur \tabularnewline
Vidéo & Chaîne officielle Apple Developer (vidéos des WWDC) & Les explications de référence des nouveautés chaque année \tabularnewline
Design & Guide des interfaces humaines d'Apple (HIG) & La bible officielle pour créer des interfaces au « goût d'Apple » \tabularnewline
Complément & La septième partie « Pratique de la programmation par IA » (chapitres 37 à 47) & La méthode complète pour développer iOS avec l'aide de l'IA \tabularnewline
\hline
\end{tabularx}
\end{center}

\begin{center}\small\color{gray}Liste des ressources pour apprendre le développement iOS (août 2026)\end{center}

\booksection{V. Conclusion : de la première app à l'énième app}

Avec le recul, tu es passé de « je ne sais pas ce qu'est SwiftUI » à : installer Xcode, apprendre Swift, dessiner des interfaces avec SwiftUI, brancher une base de données et le réseau sur ton app, puis tester, empaqueter, publier — sur ce chemin, tu as déjà franchi le tronçon le plus difficile, celui qui mène de zéro à un.

Le monde de la technologie va très vite : les frameworks se mettent à niveau, les outils changent de génération, l'IA prend en charge de plus en plus de travail de « frappe au clavier ». Mais une chose ne se démode pas : ta capacité à décomposer une idée floue en étapes claires, puis à la réaliser pas à pas. Les outils sont des amplificateurs, la réflexion est le socle. Un amplificateur peut se remplacer, le socle ne doit pas bouger.

Tu n'es plus juste "quelqu'un qui sait se servir de son téléphone" : tu peux transformer les idées de ta tête en apps qui tournent réellement sur des centaines de millions d'iPhone. À partir d'aujourd'hui, crée avec ton clavier – c'est ça, la poésie du développeur iOS.

\booksection{Références}

\begin{enumerate}
  \item Apple Inc. (2026). SwiftUI Framework Documentation. developer.apple.com.
  \item Apple Inc. (2026). Swift Programming Language Guide. swift.org.
  \item Apple Inc. (2026). SwiftData Framework Documentation. developer.apple.com.
  \item Apple Inc. (2026). App Store Review Guidelines. developer.apple.com/app-store/review.
  \item Apple Inc. (2026). Human Interface Guidelines (HIG). developer.apple.com/design.
  \item Apple Inc. (2026). Xcode 27 Release Notes and App Distribution Guide.
\end{enumerate}
